\section{Information Retrieval Benchmarks}
\label{sec:retrieval_results}

Information Retrieval (IR) tasks evaluate asymmetric search, where an incoming user query retrieves relevant document passages from a large candidate corpus. In this setting, the transformation parameters (empirical mean vector $\mu_D$, covariance matrix $\Sigma_D$, and orthogonal projection matrix $W_D$) are estimated strictly on the document corpus representations $D \in \mathbb{R}^{N \times d}$. The identical transformation parameters are applied to query vectors $Q \in \mathbb{R}^{M \times d}$:
\begin{equation}
D' = (D - \mu_D) W_D, \quad Q' = (Q - \mu_D) W_D
\end{equation}
Both transformed matrices are L2-normalized to unit length: $D' \leftarrow D' / \|D'\|_2$, $Q' \leftarrow Q' / \|Q'\|_2$. Candidate documents are ranked by inner product $S = Q' (D')^\top$, and the top 100 retrieved candidates per query are evaluated using \texttt{pytrec-eval}.

\subsection{ArguAna Counter-Argument Retrieval}
ArguAna evaluates the retrieval of counter-arguments against input argumentative claims across 8,674 candidate documents and 1,406 queries. Table~\ref{tab:arguana_ndcg}, Table~\ref{tab:arguana_mrr}, and Table~\ref{tab:arguana_recall} report nDCG@10, MRR@10, and Recall@100 across models and transformations.

% Auto-generated publication table: arguana_ndcg
\begin{table*}[t]
\centering
\small
\caption{Full-dimension evaluation on \textbf{ArguAna} (nDCG@10). Top performance across valid methods in each column is shown in bold.}
\label{tab:arguana_ndcg}
\resizebox{\textwidth}{!}{%
\begin{tabular}{l cccc cccc c}
\toprule
& \multicolumn{2}{c}{\textbf{Dedicated Models}} & \multicolumn{2}{c}{\textbf{Base LLM (Mean)}} & \multicolumn{2}{c}{\textbf{Base LLM (Last)}} & \multicolumn{3}{c}{\textbf{Tier Averages}} \\
\cmidrule(lr){2-3} \cmidrule(lr){4-5} \cmidrule(lr){6-7} \cmidrule(lr){8-10}
\textbf{Method} & \textbf{EmbGemma} & \textbf{Qwen3-Emb} & \textbf{Gemma-3} & \textbf{Qwen3.5} & \textbf{Gemma-3} & \textbf{Qwen3.5} & \textbf{Dedicated} & \textbf{Base (Mean)} & \textbf{Base (Last)} \\
\midrule
Baseline & 0.4702 & \textbf{0.4930} & 0.2938 & 0.3162 & 0.1124 & 0.1064 & \textbf{0.4816} & 0.3050 & 0.1094 \\
Standardization & 0.4754 & 0.4796 & 0.3172 & 0.3077 & 0.1299 & 0.1125 & 0.4775 & 0.3124 & 0.1212 \\
R1 (Mean Sub.) & 0.4760 & 0.4829 & 0.2988 & 0.2954 & 0.1108 & 0.1078 & 0.4794 & 0.2971 & 0.1093 \\
R2 (Subspace Def.) & 0.4697 & 0.4852 & 0.2985 & 0.2952 & 0.1121 & 0.1080 & 0.4774 & 0.2969 & 0.1100 \\
Soft-ZCA & 0.4287 & 0.4447 & \textbf{0.3599} & \textbf{0.3622} & \textbf{0.1584} & \textbf{0.1485} & 0.4367 & \textbf{0.3610} & \textbf{0.1534} \\
ABTT-1 & 0.4764 & 0.4825 & 0.3056 & 0.2986 & 0.1286 & 0.1109 & 0.4795 & 0.3021 & 0.1198 \\
ABTT-2 & \textbf{0.4767} & 0.4805 & 0.3108 & 0.3012 & 0.1338 & 0.1169 & 0.4786 & 0.3060 & 0.1253 \\
ABTT-3 & 0.4686 & 0.4792 & 0.3104 & 0.3157 & 0.1433 & 0.1301 & 0.4739 & 0.3130 & 0.1367 \\
\midrule
Rand (Control) & 0.4700 & 0.4931 & 0.2938 & 0.3163 & 0.1130 & 0.1063 & 0.4815 & 0.3051 & 0.1097 \\
MC (Control) & 0.4778 & 0.4592 & 0.2299 & 0.1924 & 0.0939 & 0.0819 & 0.4685 & 0.2111 & 0.0879 \\
\bottomrule
\end{tabular}%
}
\end{table*}

% Auto-generated publication table: arguana_mrr
\begin{table*}[t]
\centering
\small
\caption{Full-dimension evaluation on \textbf{ArguAna} (MRR@10). Top performance across valid methods in each column is shown in bold.}
\label{tab:arguana_mrr}
\resizebox{\textwidth}{!}{%
\begin{tabular}{l cccc cccc c}
\toprule
& \multicolumn{2}{c}{\textbf{Dedicated Models}} & \multicolumn{2}{c}{\textbf{Base LLM (Mean)}} & \multicolumn{2}{c}{\textbf{Base LLM (Last)}} & \multicolumn{3}{c}{\textbf{Tier Averages}} \\
\cmidrule(lr){2-3} \cmidrule(lr){4-5} \cmidrule(lr){6-7} \cmidrule(lr){8-10}
\textbf{Method} & \textbf{EmbGemma} & \textbf{Qwen3-Emb} & \textbf{Gemma-3} & \textbf{Qwen3.5} & \textbf{Gemma-3} & \textbf{Qwen3.5} & \textbf{Dedicated} & \textbf{Base (Mean)} & \textbf{Base (Last)} \\
\midrule
Baseline & 0.3311 & \textbf{0.3505} & 0.1908 & 0.2078 & 0.0723 & 0.0699 & \textbf{0.3408} & 0.1993 & 0.0711 \\
Standardization & 0.3354 & 0.3383 & 0.2024 & 0.2004 & 0.0836 & 0.0741 & 0.3369 & 0.2014 & 0.0788 \\
R1 (Mean Sub.) & 0.3364 & 0.3412 & 0.1910 & 0.1929 & 0.0704 & 0.0703 & 0.3388 & 0.1919 & 0.0704 \\
R2 (Subspace Def.) & 0.3302 & 0.3431 & 0.1905 & 0.1928 & 0.0712 & 0.0704 & 0.3366 & 0.1917 & 0.0708 \\
Soft-ZCA & 0.2969 & 0.3078 & \textbf{0.2346} & \textbf{0.2366} & \textbf{0.1002} & \textbf{0.0960} & 0.3023 & \textbf{0.2356} & \textbf{0.0981} \\
ABTT-1 & 0.3362 & 0.3410 & 0.1945 & 0.1941 & 0.0823 & 0.0731 & 0.3386 & 0.1943 & 0.0777 \\
ABTT-2 & \textbf{0.3365} & 0.3388 & 0.1982 & 0.1967 & 0.0860 & 0.0763 & 0.3377 & 0.1975 & 0.0812 \\
ABTT-3 & 0.3283 & 0.3385 & 0.1982 & 0.2053 & 0.0911 & 0.0829 & 0.3334 & 0.2017 & 0.0870 \\
\midrule
Rand (Control) & 0.3308 & 0.3506 & 0.1909 & 0.2078 & 0.0727 & 0.0700 & 0.3407 & 0.1993 & 0.0713 \\
MC (Control) & 0.3379 & 0.3160 & 0.1400 & 0.1164 & 0.0599 & 0.0542 & 0.3269 & 0.1282 & 0.0571 \\
\bottomrule
\end{tabular}%
}
\end{table*}

% Auto-generated publication table: arguana_recall
\begin{table*}[t]
\centering
\small
\caption{Full-dimension evaluation on \textbf{ArguAna} (Recall@100). Top performance across valid methods in each column is shown in bold.}
\label{tab:arguana_recall}
\resizebox{\textwidth}{!}{%
\begin{tabular}{l cccc cccc c}
\toprule
& \multicolumn{2}{c}{\textbf{Dedicated Models}} & \multicolumn{2}{c}{\textbf{Base LLM (Mean)}} & \multicolumn{2}{c}{\textbf{Base LLM (Last)}} & \multicolumn{3}{c}{\textbf{Tier Averages}} \\
\cmidrule(lr){2-3} \cmidrule(lr){4-5} \cmidrule(lr){6-7} \cmidrule(lr){8-10}
\textbf{Method} & \textbf{EmbGemma} & \textbf{Qwen3-Emb} & \textbf{Gemma-3} & \textbf{Qwen3.5} & \textbf{Gemma-3} & \textbf{Qwen3.5} & \textbf{Dedicated} & \textbf{Base (Mean)} & \textbf{Base (Last)} \\
\midrule
Baseline & 0.9936 & \textbf{0.9957} & 0.9125 & 0.9403 & 0.3890 & 0.3549 & \textbf{0.9947} & 0.9264 & 0.3720 \\
Standardization & \textbf{0.9943} & 0.9943 & 0.9630 & 0.9403 & 0.4253 & 0.3734 & 0.9943 & 0.9516 & 0.3994 \\
R1 (Mean Sub.) & \textbf{0.9943} & 0.9943 & 0.9317 & 0.9289 & 0.3940 & 0.3634 & 0.9943 & 0.9303 & 0.3787 \\
R2 (Subspace Def.) & 0.9936 & 0.9936 & 0.9324 & 0.9303 & 0.3947 & 0.3649 & 0.9936 & 0.9314 & 0.3798 \\
Soft-ZCA & 0.9723 & 0.9780 & \textbf{0.9780} & \textbf{0.9723} & \textbf{0.6394} & \textbf{0.5839} & 0.9751 & \textbf{0.9751} & \textbf{0.6117} \\
ABTT-1 & 0.9936 & 0.9936 & 0.9516 & 0.9374 & 0.4509 & 0.3947 & 0.9936 & 0.9445 & 0.4228 \\
ABTT-2 & 0.9936 & 0.9929 & 0.9516 & 0.9360 & 0.5028 & 0.4353 & 0.9932 & 0.9438 & 0.4691 \\
ABTT-3 & 0.9915 & 0.9950 & 0.9559 & 0.9502 & 0.5533 & 0.5071 & 0.9932 & 0.9531 & 0.5302 \\
\midrule
Rand (Control) & 0.9936 & 0.9957 & 0.9118 & 0.9403 & 0.3883 & 0.3535 & 0.9947 & 0.9260 & 0.3709 \\
MC (Control) & 0.9943 & 0.9929 & 0.9118 & 0.8777 & 0.3798 & 0.3450 & 0.9936 & 0.8947 & 0.3624 \\
\bottomrule
\end{tabular}%
}
\end{table*}


\paragraph{Analysis of ArguAna Results}
\begin{itemize}
\item Dedicated Embedding Models: Dedicated contrastive encoders achieve high baseline performance (mean nDCG@10 of 0.4816, MRR@10 of 0.3408, Recall@100 of 0.9947). Soft-ZCA whitening reduces dedicated nDCG@10 to 0.4367 (-0.0449) because contrastive training has already optimized the relative component weighting. Subspace deflations (R1: 0.4794; R2: 0.4774; ABTT-1: 0.4795) preserve baseline scores without metric distortion.
\item Base Causal Models under Mean Pooling: Baseline mean representations achieve 0.3050 nDCG@10, 0.1993 MRR@10, and 0.9264 Recall@100. Soft-ZCA produces the strongest recovery across all methods, raising nDCG@10 to 0.3610 (+0.0560 absolute, +18.4\% relative), MRR@10 to 0.2356 (+18.2\%), and Recall@100 to 0.9751, closing 31.7\% of the initial gap to dedicated models. ABTT-3 reaches 0.3130, and standardization reaches 0.3124.
\item Base Causal Models under Last Token Pooling: Last token representations suffer an initial penalty of 0.1956 nDCG points relative to mean pooling (baseline 0.1094 vs 0.3050). Soft-ZCA lifts last token nDCG@10 to 0.1534 (+40.2\%) and Recall@100 from 0.3720 to 0.6117 (+64.4\%).
\item Negative Controls: Random unit direction deflation (\texttt{rand}) yields 0.4815 for dedicated and 0.3051 for base mean pooling, demonstrating zero statistical shift from baseline (0.4816 and 0.3050). Mean centering without L2 renormalization (\texttt{mc}) causes base mean pooling nDCG@10 to collapse from 0.3050 to 0.2111 (-30.8\%), confirming that hyperspherical projection is mathematically essential.
\end{itemize}


\subsection{FiQA2018 Financial Question Answering}
FiQA2018 evaluates financial opinion question answering across 57,638 document passages and 648 queries. Table~\ref{tab:fiqa_ndcg}, Table~\ref{tab:fiqa_mrr}, and Table~\ref{tab:fiqa_recall} report nDCG@10, MRR@10, and Recall@100.

% Auto-generated publication table: fiqa_ndcg
\begin{table*}[t]
\centering
\small
\caption{Full-dimension evaluation on \textbf{FiQA2018} (nDCG@10). Top performance across valid methods in each column is shown in bold.}
\label{tab:fiqa_ndcg}
\resizebox{\textwidth}{!}{%
\begin{tabular}{l cccc cccc c}
\toprule
& \multicolumn{2}{c}{\textbf{Dedicated Models}} & \multicolumn{2}{c}{\textbf{Base LLM (Mean)}} & \multicolumn{2}{c}{\textbf{Base LLM (Last)}} & \multicolumn{3}{c}{\textbf{Tier Averages}} \\
\cmidrule(lr){2-3} \cmidrule(lr){4-5} \cmidrule(lr){6-7} \cmidrule(lr){8-10}
\textbf{Method} & \textbf{EmbGemma} & \textbf{Qwen3-Emb} & \textbf{Gemma-3} & \textbf{Qwen3.5} & \textbf{Gemma-3} & \textbf{Qwen3.5} & \textbf{Dedicated} & \textbf{Base (Mean)} & \textbf{Base (Last)} \\
\midrule
Baseline & 0.4739 & 0.4737 & 0.0589 & 0.0419 & 0.0246 & 0.0026 & 0.4738 & 0.0504 & 0.0136 \\
Standardization & 0.4704 & 0.4656 & 0.0681 & 0.0385 & 0.0187 & 0.0019 & 0.4680 & 0.0533 & 0.0103 \\
R1 (Mean Sub.) & 0.4731 & 0.4725 & 0.0564 & 0.0420 & 0.0225 & 0.0025 & 0.4728 & 0.0492 & 0.0125 \\
R2 (Subspace Def.) & \textbf{0.4765} & 0.4730 & 0.0572 & 0.0433 & 0.0223 & 0.0028 & 0.4748 & 0.0502 & 0.0126 \\
Soft-ZCA & 0.4056 & 0.4151 & \textbf{0.0968} & \textbf{0.0615} & \textbf{0.0389} & \textbf{0.0042} & 0.4103 & \textbf{0.0791} & \textbf{0.0215} \\
ABTT-1 & 0.4736 & \textbf{0.4769} & 0.0568 & 0.0378 & 0.0222 & 0.0026 & \textbf{0.4752} & 0.0473 & 0.0124 \\
ABTT-2 & 0.4641 & 0.4765 & 0.0555 & 0.0422 & 0.0272 & 0.0036 & 0.4703 & 0.0489 & 0.0154 \\
ABTT-3 & 0.4625 & 0.4676 & 0.0689 & 0.0446 & 0.0254 & 0.0036 & 0.4651 & 0.0567 & 0.0145 \\
\midrule
Rand (Control) & 0.4736 & 0.4726 & 0.0587 & 0.0416 & 0.0247 & 0.0026 & 0.4731 & 0.0502 & 0.0136 \\
MC (Control) & 0.4539 & 0.4597 & 0.0181 & 0.0126 & 0.0200 & 0.0016 & 0.4568 & 0.0154 & 0.0108 \\
\bottomrule
\end{tabular}%
}
\end{table*}

% Auto-generated publication table: fiqa_mrr
\begin{table*}[t]
\centering
\small
\caption{Full-dimension evaluation on \textbf{FiQA2018} (MRR@10). Top performance across valid methods in each column is shown in bold.}
\label{tab:fiqa_mrr}
\resizebox{\textwidth}{!}{%
\begin{tabular}{l cccc cccc c}
\toprule
& \multicolumn{2}{c}{\textbf{Dedicated Models}} & \multicolumn{2}{c}{\textbf{Base LLM (Mean)}} & \multicolumn{2}{c}{\textbf{Base LLM (Last)}} & \multicolumn{3}{c}{\textbf{Tier Averages}} \\
\cmidrule(lr){2-3} \cmidrule(lr){4-5} \cmidrule(lr){6-7} \cmidrule(lr){8-10}
\textbf{Method} & \textbf{EmbGemma} & \textbf{Qwen3-Emb} & \textbf{Gemma-3} & \textbf{Qwen3.5} & \textbf{Gemma-3} & \textbf{Qwen3.5} & \textbf{Dedicated} & \textbf{Base (Mean)} & \textbf{Base (Last)} \\
\midrule
Baseline & 0.5467 & 0.5509 & 0.0700 & 0.0567 & 0.0299 & 0.0034 & 0.5488 & 0.0634 & 0.0167 \\
Standardization & 0.5404 & 0.5440 & 0.0868 & 0.0486 & 0.0258 & 0.0027 & 0.5422 & 0.0677 & 0.0142 \\
R1 (Mean Sub.) & 0.5424 & 0.5509 & 0.0666 & 0.0541 & 0.0284 & 0.0035 & 0.5466 & 0.0604 & 0.0159 \\
R2 (Subspace Def.) & \textbf{0.5481} & 0.5496 & 0.0686 & 0.0568 & 0.0296 & 0.0036 & \textbf{0.5489} & 0.0627 & 0.0166 \\
Soft-ZCA & 0.4823 & 0.4929 & \textbf{0.1132} & \textbf{0.0687} & \textbf{0.0449} & \textbf{0.0041} & 0.4876 & \textbf{0.0909} & \textbf{0.0245} \\
ABTT-1 & 0.5442 & 0.5536 & 0.0697 & 0.0475 & 0.0268 & 0.0035 & \textbf{0.5489} & 0.0586 & 0.0152 \\
ABTT-2 & 0.5359 & \textbf{0.5558} & 0.0682 & 0.0537 & 0.0322 & 0.0040 & 0.5458 & 0.0609 & 0.0181 \\
ABTT-3 & 0.5340 & 0.5440 & 0.0855 & 0.0566 & 0.0302 & 0.0037 & 0.5390 & 0.0710 & 0.0170 \\
\midrule
Rand (Control) & 0.5460 & 0.5502 & 0.0699 & 0.0566 & 0.0300 & 0.0034 & 0.5481 & 0.0633 & 0.0167 \\
MC (Control) & 0.5245 & 0.5360 & 0.0252 & 0.0164 & 0.0264 & 0.0019 & 0.5302 & 0.0208 & 0.0142 \\
\bottomrule
\end{tabular}%
}
\end{table*}

% Auto-generated publication table: fiqa_recall
\begin{table*}[t]
\centering
\small
\caption{Full-dimension evaluation on \textbf{FiQA2018} (Recall@100). Top performance across valid methods in each column is shown in bold.}
\label{tab:fiqa_recall}
\resizebox{\textwidth}{!}{%
\begin{tabular}{l cccc cccc c}
\toprule
& \multicolumn{2}{c}{\textbf{Dedicated Models}} & \multicolumn{2}{c}{\textbf{Base LLM (Mean)}} & \multicolumn{2}{c}{\textbf{Base LLM (Last)}} & \multicolumn{3}{c}{\textbf{Tier Averages}} \\
\cmidrule(lr){2-3} \cmidrule(lr){4-5} \cmidrule(lr){6-7} \cmidrule(lr){8-10}
\textbf{Method} & \textbf{EmbGemma} & \textbf{Qwen3-Emb} & \textbf{Gemma-3} & \textbf{Qwen3.5} & \textbf{Gemma-3} & \textbf{Qwen3.5} & \textbf{Dedicated} & \textbf{Base (Mean)} & \textbf{Base (Last)} \\
\midrule
Baseline & 0.7959 & \textbf{0.8154} & 0.2443 & 0.1717 & 0.1169 & 0.0206 & \textbf{0.8056} & 0.2080 & 0.0688 \\
Standardization & 0.7921 & 0.8018 & 0.2570 & 0.1611 & 0.1172 & 0.0260 & 0.7970 & 0.2091 & 0.0716 \\
R1 (Mean Sub.) & 0.7917 & 0.8071 & 0.2208 & 0.1775 & 0.1148 & 0.0198 & 0.7994 & 0.1991 & 0.0673 \\
R2 (Subspace Def.) & \textbf{0.7981} & 0.8104 & 0.2352 & 0.1746 & 0.1115 & 0.0236 & 0.8042 & 0.2049 & 0.0675 \\
Soft-ZCA & 0.7137 & 0.7400 & \textbf{0.3663} & \textbf{0.2692} & \textbf{0.1898} & \textbf{0.0644} & 0.7268 & \textbf{0.3177} & \textbf{0.1271} \\
ABTT-1 & 0.7966 & 0.8077 & 0.2399 & 0.1608 & 0.1232 & 0.0265 & 0.8021 & 0.2003 & 0.0749 \\
ABTT-2 & 0.7923 & 0.7985 & 0.2393 & 0.1774 & 0.1451 & 0.0470 & 0.7954 & 0.2084 & 0.0960 \\
ABTT-3 & 0.7881 & 0.7953 & 0.2725 & 0.1822 & 0.1377 & 0.0390 & 0.7917 & 0.2274 & 0.0883 \\
\midrule
Rand (Control) & 0.7959 & 0.8154 & 0.2438 & 0.1717 & 0.1177 & 0.0206 & 0.8056 & 0.2077 & 0.0692 \\
MC (Control) & 0.7809 & 0.8028 & 0.1026 & 0.0943 & 0.0965 & 0.0200 & 0.7919 & 0.0984 & 0.0582 \\
\bottomrule
\end{tabular}%
}
\end{table*}


\paragraph{Analysis of FiQA2018 Results}
\begin{itemize}
\item Initial Disparity: Raw base causal models perform poorly on specialized financial QA, achieving an average baseline nDCG@10 of 0.0504 under mean pooling and 0.0136 under last token pooling, trailing dedicated models (0.4738) by 0.4234 points.
\item Post-Processing Gains on Base Models: Soft-ZCA delivers a +56.9\% relative gain on base mean pooling, lifting nDCG@10 from 0.0504 to 0.0791, MRR@10 from 0.0634 to 0.0909 (+43.4\%), and Recall@100 from 0.2080 to 0.3177 (+52.7\%). ABTT-3 reaches 0.0567 nDCG@10 and 0.2274 Recall@100.
\item Extreme Collapse on Last Token Pooling: Qwen3.5-0.8B-Base under last token pooling retrieves almost no relevant financial passages (0.0026 nDCG@10, 0.0034 MRR@10, and 0.0206 Recall@100). Soft-ZCA raises recall to 0.0644, but last token pooling remains deeply compromised by autoregressive token specialization.
\item Dedicated Model Stability: Dedicated models preserve peak performance under R2 (0.4748 nDCG vs 0.4738 baseline) and ABTT-1 (0.4752 nDCG), whereas Soft-ZCA causes an over-smoothing drop to 0.4103.
\end{itemize}


\subsection{SCIDOCS Scientific Citation Retrieval}
SCIDOCS evaluates paper citation matching across 25,657 paper abstracts and 1,000 title queries with 29,928 relevant citation links. Table~\ref{tab:scidocs_ndcg}, Table~\ref{tab:scidocs_mrr}, and Table~\ref{tab:scidocs_recall} report nDCG@10, MRR@10, and Recall@100.

% Auto-generated publication table: scidocs_ndcg
\begin{table*}[t]
\centering
\small
\caption{Full-dimension evaluation on \textbf{SCIDOCS} (nDCG@10). Top performance across valid methods in each column is shown in bold.}
\label{tab:scidocs_ndcg}
\resizebox{\textwidth}{!}{%
\begin{tabular}{l cccc cccc c}
\toprule
& \multicolumn{2}{c}{\textbf{Dedicated Models}} & \multicolumn{2}{c}{\textbf{Base LLM (Mean)}} & \multicolumn{2}{c}{\textbf{Base LLM (Last)}} & \multicolumn{3}{c}{\textbf{Tier Averages}} \\
\cmidrule(lr){2-3} \cmidrule(lr){4-5} \cmidrule(lr){6-7} \cmidrule(lr){8-10}
\textbf{Method} & \textbf{EmbGemma} & \textbf{Qwen3-Emb} & \textbf{Gemma-3} & \textbf{Qwen3.5} & \textbf{Gemma-3} & \textbf{Qwen3.5} & \textbf{Dedicated} & \textbf{Base (Mean)} & \textbf{Base (Last)} \\
\midrule
Baseline & 0.1948 & \textbf{0.2439} & 0.0302 & 0.0268 & 0.0057 & 0.0042 & 0.2194 & 0.0285 & 0.0050 \\
Standardization & 0.1940 & 0.2402 & 0.0420 & 0.0193 & 0.0062 & 0.0049 & 0.2171 & 0.0307 & 0.0055 \\
R1 (Mean Sub.) & 0.1992 & 0.2403 & 0.0222 & 0.0216 & 0.0052 & 0.0048 & 0.2198 & 0.0219 & 0.0050 \\
R2 (Subspace Def.) & 0.1966 & 0.2405 & 0.0244 & 0.0229 & 0.0055 & 0.0048 & 0.2185 & 0.0237 & 0.0051 \\
Soft-ZCA & 0.1686 & 0.2161 & \textbf{0.0679} & \textbf{0.0536} & \textbf{0.0106} & \textbf{0.0080} & 0.1924 & \textbf{0.0607} & \textbf{0.0093} \\
ABTT-1 & 0.1994 & 0.2406 & 0.0226 & 0.0197 & 0.0053 & 0.0042 & 0.2200 & 0.0212 & 0.0048 \\
ABTT-2 & \textbf{0.2012} & 0.2405 & 0.0214 & 0.0189 & 0.0057 & 0.0044 & \textbf{0.2208} & 0.0201 & 0.0051 \\
ABTT-3 & 0.1996 & 0.2380 & 0.0226 & 0.0172 & 0.0068 & 0.0055 & 0.2188 & 0.0199 & 0.0061 \\
\midrule
Rand (Control) & 0.1950 & 0.2440 & 0.0299 & 0.0267 & 0.0058 & 0.0042 & 0.2195 & 0.0283 & 0.0050 \\
MC (Control) & 0.1930 & 0.2288 & 0.0093 & 0.0111 & 0.0061 & 0.0053 & 0.2109 & 0.0102 & 0.0057 \\
\bottomrule
\end{tabular}%
}
\end{table*}

% Auto-generated publication table: scidocs_mrr
\begin{table*}[t]
\centering
\small
\caption{Full-dimension evaluation on \textbf{SCIDOCS} (MRR@10). Top performance across valid methods in each column is shown in bold.}
\label{tab:scidocs_mrr}
\resizebox{\textwidth}{!}{%
\begin{tabular}{l cccc cccc c}
\toprule
& \multicolumn{2}{c}{\textbf{Dedicated Models}} & \multicolumn{2}{c}{\textbf{Base LLM (Mean)}} & \multicolumn{2}{c}{\textbf{Base LLM (Last)}} & \multicolumn{3}{c}{\textbf{Tier Averages}} \\
\cmidrule(lr){2-3} \cmidrule(lr){4-5} \cmidrule(lr){6-7} \cmidrule(lr){8-10}
\textbf{Method} & \textbf{EmbGemma} & \textbf{Qwen3-Emb} & \textbf{Gemma-3} & \textbf{Qwen3.5} & \textbf{Gemma-3} & \textbf{Qwen3.5} & \textbf{Dedicated} & \textbf{Base (Mean)} & \textbf{Base (Last)} \\
\midrule
Baseline & 0.3345 & \textbf{0.4072} & 0.0551 & 0.0456 & 0.0114 & 0.0085 & 0.3708 & 0.0503 & 0.0100 \\
Standardization & 0.3325 & 0.4025 & 0.0723 & 0.0340 & 0.0139 & 0.0097 & 0.3675 & 0.0531 & 0.0118 \\
R1 (Mean Sub.) & 0.3404 & 0.4007 & 0.0404 & 0.0361 & 0.0109 & 0.0097 & 0.3706 & 0.0383 & 0.0103 \\
R2 (Subspace Def.) & 0.3391 & 0.4006 & 0.0448 & 0.0377 & 0.0110 & 0.0095 & 0.3699 & 0.0413 & 0.0103 \\
Soft-ZCA & 0.3006 & 0.3696 & \textbf{0.1176} & \textbf{0.0909} & \textbf{0.0200} & \textbf{0.0169} & 0.3351 & \textbf{0.1043} & \textbf{0.0185} \\
ABTT-1 & 0.3397 & 0.4021 & 0.0414 & 0.0337 & 0.0107 & 0.0085 & 0.3709 & 0.0376 & 0.0096 \\
ABTT-2 & \textbf{0.3431} & 0.4005 & 0.0388 & 0.0324 & 0.0120 & 0.0086 & \textbf{0.3718} & 0.0356 & 0.0103 \\
ABTT-3 & 0.3409 & 0.3956 & 0.0414 & 0.0311 & 0.0140 & 0.0111 & 0.3682 & 0.0362 & 0.0126 \\
\midrule
Rand (Control) & 0.3346 & 0.4073 & 0.0546 & 0.0455 & 0.0116 & 0.0085 & 0.3710 & 0.0501 & 0.0100 \\
MC (Control) & 0.3294 & 0.3814 & 0.0171 & 0.0234 & 0.0133 & 0.0127 & 0.3554 & 0.0203 & 0.0130 \\
\bottomrule
\end{tabular}%
}
\end{table*}

% Auto-generated publication table: scidocs_recall
\begin{table*}[t]
\centering
\small
\caption{Full-dimension evaluation on \textbf{SCIDOCS} (Recall@100). Top performance across valid methods in each column is shown in bold.}
\label{tab:scidocs_recall}
\resizebox{\textwidth}{!}{%
\begin{tabular}{l cccc cccc c}
\toprule
& \multicolumn{2}{c}{\textbf{Dedicated Models}} & \multicolumn{2}{c}{\textbf{Base LLM (Mean)}} & \multicolumn{2}{c}{\textbf{Base LLM (Last)}} & \multicolumn{3}{c}{\textbf{Tier Averages}} \\
\cmidrule(lr){2-3} \cmidrule(lr){4-5} \cmidrule(lr){6-7} \cmidrule(lr){8-10}
\textbf{Method} & \textbf{EmbGemma} & \textbf{Qwen3-Emb} & \textbf{Gemma-3} & \textbf{Qwen3.5} & \textbf{Gemma-3} & \textbf{Qwen3.5} & \textbf{Dedicated} & \textbf{Base (Mean)} & \textbf{Base (Last)} \\
\midrule
Baseline & 0.4230 & \textbf{0.5292} & 0.1108 & 0.1165 & 0.0366 & 0.0270 & 0.4761 & 0.1136 & 0.0318 \\
Standardization & 0.4331 & 0.5190 & 0.1726 & 0.0855 & 0.0329 & 0.0243 & 0.4761 & 0.1291 & 0.0286 \\
R1 (Mean Sub.) & \textbf{0.4424} & 0.5217 & 0.0899 & 0.0957 & 0.0293 & 0.0243 & \textbf{0.4821} & 0.0928 & 0.0268 \\
R2 (Subspace Def.) & 0.4396 & 0.5193 & 0.1002 & 0.1006 & 0.0311 & 0.0247 & 0.4794 & 0.1004 & 0.0279 \\
Soft-ZCA & 0.3566 & 0.4519 & \textbf{0.2338} & \textbf{0.2172} & \textbf{0.0914} & \textbf{0.0744} & 0.4043 & \textbf{0.2255} & \textbf{0.0829} \\
ABTT-1 & 0.4394 & 0.5171 & 0.0880 & 0.0865 & 0.0331 & 0.0271 & 0.4782 & 0.0872 & 0.0301 \\
ABTT-2 & 0.4412 & 0.5190 & 0.0838 & 0.0785 & 0.0390 & 0.0339 & 0.4801 & 0.0812 & 0.0365 \\
ABTT-3 & 0.4374 & 0.5152 & 0.0895 & 0.0715 & 0.0499 & 0.0478 & 0.4763 & 0.0805 & 0.0489 \\
\midrule
Rand (Control) & 0.4216 & 0.5292 & 0.1104 & 0.1165 & 0.0366 & 0.0270 & 0.4754 & 0.1134 & 0.0318 \\
MC (Control) & 0.4439 & 0.5124 & 0.0474 & 0.0622 & 0.0278 & 0.0247 & 0.4782 & 0.0548 & 0.0263 \\
\bottomrule
\end{tabular}%
}
\end{table*}


\paragraph{Analysis of SCIDOCS Results}
\begin{itemize}
\item Base Representation Deficits: Raw base causal models achieve only 0.0285 nDCG@10 and 0.0503 MRR@10 under mean pooling, and 0.0050 nDCG@10 under last token pooling, compared to 0.2194 and 0.3708 for dedicated encoders.
\item Soft-ZCA Recovery: Soft-ZCA more than doubles base model retrieval quality, lifting nDCG@10 from 0.0285 to 0.0607 (+113.0\% relative increase), MRR@10 from 0.0503 to 0.1043 (+107.4\%), and Recall@100 from 0.1136 to 0.2255 (+98.5\%), closing 16.9\% of the initial gap to dedicated models.
\item Failure of 1D Deflation: Unlike Soft-ZCA, simple mean subtraction (R1: 0.0219; R2: 0.0237) and coordinate standardization (0.0307) provide little to no gain on SCIDOCS. In scientific abstracts, representational crowding is driven by higher-order covariance coupling rather than origin shift alone.
\item Dedicated Encoders Under Deflation: Dedicated models slightly improve under light deflation (ABTT-2 reaches 0.2208 nDCG vs 0.2194 baseline; R1 reaches 0.4821 Recall vs 0.4761 baseline), demonstrating that removing the top two principal background directions purges residual citation dataset noise.
\end{itemize}
